\chapter{Selected Implementation Details}
\label{app:code}
\section{Source and Reproducibility}
All code excerpts below are taken from the final-experiment branch of the repository. They are illustrative, not independent implementations. The authoritative code is the frozen repository commit recorded in the experiment manifest.
\section{Single-Response API}
\begin{lstlisting}[language=Python,caption={Single-response entry point in src/cultverify/pipeline.py.}]
def verify(self, prompt: str, response: str) -> CandidateResult:
    self._inputs(prompt, response)
    session = SemanticSession(self.llm, self.config)
    try:
        context, plan = plan_prompt(session, prompt, self.rubric)
        planning_error = None
    except StageError as exc:
        context, plan = ContextFrame(), DimensionPlan(
            dimensions=(), reasoning="Planning failed"
        )
        planning_error = str(exc)
    return self._verify(
        prompt, response, context, plan,
        tuple(session.calls), planning_error
    )
\end{lstlisting}
The entry point first plans the prompt and then evaluates the response. It also handles a planning failure by passing a recorded error into the verification routine. The final experiment uses this API, not the four-response ranking API.
\section{Exact Response Spans}
\begin{lstlisting}[language=Python,caption={Response span selection in src/cultverify/targets.py.}]
def response_spans(response):
    spans = []
    for line in response.splitlines():
        line = line.strip()
        if not line:
            continue
        pieces = tuple(
            piece.strip()
            for piece in _SENTENCE_BOUNDARY.split(line)
            if piece.strip()
        )
        for piece in pieces or (line,):
            spans.append({
                "span_id": f"S{len(spans) + 1:03d}",
                "text": piece,
            })
    if response and not spans:
        spans.append({"span_id": "S001", "text": response})
    return tuple(spans)
\end{lstlisting}
This function gives the model stable span IDs to choose from. Python retains the actual response text, so the model is not required to reproduce arbitrary quotations exactly. This is important for later checking which words were actually evaluated.
\section{Evidence Sufficiency Rule}
\begin{lstlisting}[language=Python,caption={Insufficient-evidence guard in src/cultverify/scoring.py.}]
def compare_target(session, target, memo):
    if memo.sufficiency == "insufficient":
        return TargetVerdict(
            target_id=target.target_id,
            memo_id=memo.memo_id,
            verdict="insufficient",
            reasoning=(
                "Frozen evidence memo is insufficient; "
                "no directional verdict is permitted."
            ),
        )
\end{lstlisting}
The excerpt shows the beginning of the actual function. The remaining branch calls the target-comparator model when evidence is sufficient. This guard prevents an insufficient memo from being turned into an unsupported directional verdict.
\section{Dataset Hashing}
\begin{lstlisting}[language=Python,caption={Canonical semantic hash from scripts/run_final_experiment.py.}]
payload = [
    {
        "generator_model": row["generator_model"],
        "item_id": row["item_id"],
        "max_tokens": row["max_tokens"],
        "prompt_sha256": row["prompt_sha256"],
        "response_sha256": row["response_sha256"],
        "temperature": row["temperature"],
        "top_p": row["top_p"],
    }
    for row in rows
]
canonical = json.dumps(
    payload, ensure_ascii=False, sort_keys=True,
    separators=(",", ":")
)
return hashlib.sha256(canonical.encode("utf-8")).hexdigest()
\end{lstlisting}
This freezes the ordered semantic inputs without depending on how a CSV writer formats newlines or quotation marks. Per-row prompt and response hashes are checked separately before this digest is accepted.
\section{Limits of Code Inspection}
These snippets demonstrate explicit implementation choices, not the empirical success of those choices. Cultural judgment quality must be assessed from final results and trace review.